\section{A genuine finite arithmetic coefficient model of the same affine equations}

\begin{definition}[The actual finite-type integer coefficient ring contains every specified coefficient]
\label{def:arithmetic-finite-coefficient-ring}
\lean{EnvelopeProperty.ArithmeticCoefficientModel.coefficientRing,EnvelopeProperty.ArithmeticCoefficientModel.extra_mem_coefficientRing,EnvelopeProperty.ArithmeticCoefficientModel.coeff_mem_coefficientRing,EnvelopeProperty.ArithmeticCoefficientModel.coefficientRing_finiteType,EnvelopeProperty.ArithmeticCoefficientModel.coefficientRing_arithmetic}
\leanok
Let $k$ be the original characteristic-zero field, let
$f_i\in k[x_j\mid j\in\sigma]$ be an original finite family of polynomial
equations, and let $E\subset k$ be any specified finite set of extra coefficients.
Define the ACTUAL subalgebra of the SAME original field
\[
 R=\mathbb Z[\,\operatorname{coeff}(f_i)\ (i),\ E\,]\subset k.
\]
Then $R$ is a finite-type $\mathbb Z$-algebra, an integral domain of
characteristic zero. It contains every original polynomial coefficient and
every specified element of $E$, including zero coefficients.
No smoothness of $R$ over $\mathbb Z$ is asserted or assumed.
\end{definition}
\begin{proof}
Take the actual union of the finite coefficient supports of the finitely
many original polynomials and the specified finite set $E$, and form its
actual $\mathbb Z$-algebra closure inside $k$. The union is finite, so this
closure is generated by finitely many actual elements and is of finite type.
Every listed coefficient and extra element lies in the closure by construction;
zero lies in every subalgebra. Its original inclusion into the original field
is injective, giving a domain and preserving the actual integer map, hence
characteristic zero. This constructs a subring, not a subfield containing
all inverses of integers, so it is a genuine arithmetic coefficient base.
\end{proof}

\begin{theorem}[A same-presentation literal arithmetic model and its whole scalar-extension isomorphism]
\label{thm:arithmetic-same-presentation-model}
\lean{EnvelopeProperty.ArithmeticCoefficientModel.presentationRelations,EnvelopeProperty.ArithmeticCoefficientModel.map_presentationRelations,EnvelopeProperty.ArithmeticCoefficientModel.coeff_presentationRelations,EnvelopeProperty.ArithmeticCoefficientModel.presentationModel,EnvelopeProperty.ArithmeticCoefficientModel.presentationModel_finitePresentation,EnvelopeProperty.ArithmeticCoefficientModel.presentationTensorEquiv,EnvelopeProperty.ArithmeticCoefficientModel.presentationTensorEquiv_tmul,EnvelopeProperty.ArithmeticCoefficientModel.presentationTensorEquiv_coefficient,EnvelopeProperty.ArithmeticCoefficientModel.presentationTensorEquiv_variable,EnvelopeProperty.ArithmeticCoefficientModel.presentationTensorEquiv_whole_polynomial}
\leanok
\uses{def:arithmetic-finite-coefficient-ring}
Let $P$ be an actual presentation of an original $k$-algebra $A$, with
original generators $z_j\in A$ and finitely many original polynomial
relations $f_i$. Use the SAME actual $R$ from the preceding node for these
relations and any specified finite extras. Each actual original relation
has a polynomial $f_{R,i}\in R[x_j]$ whose coefficient inclusion maps it
back to exactly $f_i$, coefficient by coefficient. Put literally
\[ A_R=R[x_j]/(f_{R,i}). \]
If the original variable set is finite, $A_R$ is finitely presented over $R$.
There is a constructed ACTUAL whole $k$-algebra equivalence
\[ k\otimes_R A_R\simeq_k A \]
which, for EVERY original $c\in k$ and EVERY $p\in R[x_j]$, sends
\[
 c\otimes\overline p\longmapsto c\,p(z_j)
\]
using the SAME original coefficient inclusion. In particular all original
base coefficients and all original presentation generators are retained.
The entire polynomial/quotient diagram commutes; no base-change
isomorphism, coefficient containment or model certificate is an input.
\end{theorem}
\begin{proof}
All relation coefficients lie in the actual image of $R$ in $k$ by the
preceding node. Restrict those coefficients to $R$, giving the stated
relations. The original inclusion is injective, so these coefficient
readbacks agree with every original coefficient, including zero.
Use the native presentation construction over a ring containing its
coefficients: its model is exactly the polynomial quotient by the lifted
relations, and its native tensor-model equivalence is constructed from the
original presentation relation/kernel equality. Equivalently, scalar
extension of the polynomial algebra is the same original $k[x_j]$;
the extended relation ideal is generated by the SAME original $f_i$.
Taking the actual quotient gives precisely $A$, without a flatness or
faithfulness assumption. The native forward map on every pure tensor is
the displayed original evaluation formula. The relation family is finite,
so with finitely many variables the literal model is finitely presented.
The formula on every polynomial, and the algebra-map equality on every
original coefficient and variable, give the whole commutative diagram.
No purely general tensor/quotient framework is reproved here.
\end{proof}

\begin{theorem}[The literal original finite polynomial quotient is the same base change]
\label{thm:arithmetic-literal-polynomial-quotient-model}
\lean{EnvelopeProperty.ArithmeticCoefficientModel.polynomialPresentation,EnvelopeProperty.ArithmeticCoefficientModel.polynomialRelations,EnvelopeProperty.ArithmeticCoefficientModel.map_polynomialRelations,EnvelopeProperty.ArithmeticCoefficientModel.polynomialModel,EnvelopeProperty.ArithmeticCoefficientModel.polynomialModel_finitePresentation,EnvelopeProperty.ArithmeticCoefficientModel.polynomialTensorEquiv,EnvelopeProperty.ArithmeticCoefficientModel.polynomialTensorEquiv_tmul,EnvelopeProperty.ArithmeticCoefficientModel.polynomialTensorEquiv_whole_quotient}
\leanok
\uses{thm:arithmetic-same-presentation-model}
For the original finite equation family $f_i$ and extras $E$, retain the
SAME actual $R$, literal descended equations $f_{R,i}$ and literal model
$A_R=R[x_j]/(f_{R,i})$ above. The constructed whole equivalence is
\[ k\otimes_R A_R\simeq_k k[x_j]/(f_i). \]
Every original coefficient $c\in k$, every original variable and the ENTIRE
original quotient map are preserved. Explicitly every $c\otimes\overline p$
maps to $c\,\overline{p}$ with coefficients included in the same original
$k$, and the complete original $k[x_j]$ quotient map agrees with the
polynomial scalar-extension map followed by the model quotient and this
equivalence. The cases of no equations, no extras and no variables are included.
\end{theorem}
\begin{proof}
Apply the preceding theorem to the native actual presentation of the
literal original quotient by precisely the given family $f_i$. Its relations
are definitionally that family, and its generators are the original
quotient classes of $x_j$. The whole evaluation formula is therefore exactly
the original quotient map on every polynomial. Native polynomial scalar
extension and quotient tensor maps yield the displayed whole map square;
equality of algebra maps follows on every original variable and every
original coefficient. No replacement presentation or target isomorphism
is supplied as a hypothesis.
\end{proof}

\begin{theorem}[The original finitely presented algebra internally supplies its actual arithmetic model]
\label{thm:arithmetic-original-finite-presentation-model}
\lean{EnvelopeProperty.ArithmeticCoefficientModel.exists_original_finitePresentation_model}
\leanok
\uses{thm:arithmetic-same-presentation-model,thm:arithmetic-literal-polynomial-quotient-model}
Given the SAME original finitely presented $k$-algebra $A$ and arbitrary
specified finite extra coefficients in $k$, an actual finite presentation
of that SAME $A$ is internally obtained. The preceding construction gives
an ACTUAL finite-type integer subalgebra $R\subset k$, actual descended
finite relations, literal finitely presented model $A_R$, and the actual
whole equivalence $k\otimes_R A_R\simeq_k A$. It retains every original
coefficient and each generator of the internally produced presentation,
with the entire original polynomial quotient/evaluation map equality.
\end{theorem}
\begin{proof}
Use the native finite-presentation theorem to choose finite variable and
relation sets and an actual presentation of the SAME original $A$.
Apply the same already proved construction to that presentation and the
specified extras; its constructed tensor equivalence has target exactly
$A$. No model or target-equivalence certificate is assumed.
\end{proof}

This is the affine finite-coefficient producer for subsequent descent of
the SAME original $X$, its line bundle, and its ample witness. Finite
additional overlap, inverse, transition and embedding coefficients may be
included among the extras so that the whole displayed maps retain them.
Actual gluing, projective-family/line descent, a smooth arithmetic base,
and all vanishing statements remain separate. In particular $R\to k$
need not be faithfully flat; neither smoothness, flatness nor integrality
of $A_R$ follows here from such properties of its original generic fibre.
